% GENERATED FILE. Edit canonical lessons, then run npm run generate:books.
% canonical-source-hash: fnv1a64:1376658344f625fe
% canonical-lessons: ES-C378-selva, ES-C378-prado, ES-C378-volcan, ES-C378-cumbre, ES-C378-edificio

\chapter{Five Places You Can Ask About}
\label{ch:la-selva-y-el-prado}
\hlchaptermodality{378}

\begin{chapteropening}
\textbf{By the end of this chapter:} \emph{I can say la selva, el prado, el volcán, la cumbre and el edificio, and ask where five kinds of place are.}

Complete ES-C378-edificio: ask where five kinds of place are, and say one thing about each.
\end{chapteropening}

\section[la selva]{la selva — the wood inside the word savage}

\label{lesson:ES-C378-selva}



\begin{quote}
Say \emph{the story}, then \emph{the dance}. (\emph{La historia}; \emph{El baile}.)
\end{quote}

\subsection*{What to know first}
\begin{itemize}
\raggedright
  \item el bosque — the cooler kind.
  \item el árbol — what fills it.
\end{itemize}

\begin{sounds}
\begin{itemize}
\raggedright
  \item \emph{SEL-va}, two beats with the weight on the first. The \emph{l} is clear and forward, not the dark \emph{l} English uses at the end of a syllable. $\to$ pronunciation reference
\end{itemize}
\end{sounds}

\begin{cousinweb}
\textbf{la selva} — \textquotedblleft{}jungle.\textquotedblright{}

Latin \textbf{silva} was a wood. Spanish kept it as \emph{selva} for the thick, steaming kind. English took it three separate times, and the third is the one nobody expects.

\textbf{Sylvan} means \textquotedblleft{}of the woods\textquotedblright{} — and its \emph{y} is a mistake. Somebody decided the word must be Greek, from \emph{hýlē} \textquotedblleft{}forest,\textquotedblright{} and spelled it accordingly. The spelling stuck; the theory did not.

\textbf{Pennsylvania} is \textquotedblleft{}Penn's woods.\textquotedblright{} A detail here gets told wrong nearly every time: it does not honour William Penn the Quaker founder. Charles II named it for \textbf{his father}, Admiral Penn, and the son was reportedly uncomfortable about it.

And then \textbf{savage}. Latin \textbf{silvaticus} meant \textquotedblleft{}of the woods\textquotedblright{}; Vulgar Latin reshaped it to \emph{salvaticus}; Old French wore it down to \emph{sauvage}; English took that. So the word for wild and cruel began as the word for someone who lives among trees — a judgement passed by people who lived in towns.
\end{cousinweb}

\begin{grammarlens}[title={What you've built}]
You can now name a thick forest: \emph{la selva}.
\end{grammarlens}

\subsection*{Your turn}
Say these aloud:

\begin{itemize}
\raggedright
  \item \emph{la selva}
  \item \emph{en la selva}
  \item \emph{¿dónde está la selva?}
\end{itemize}

\begin{tcolorbox}[breakable,colback=teal!4,colframe=teal!35!black,title={Before you move on}]
Which English word for wildness comes from \emph{silvaticus}? (\emph{Savage} — \textquotedblleft{}of the woods\textquotedblright{}.) And whom is Pennsylvania named for? (Admiral Penn, the founder's father.)
\end{tcolorbox}
\section[el prado]{el prado — a museum named after a field}

\label{lesson:ES-C378-prado}



\begin{quote}
Say \emph{the dance}, then \emph{the jungle}. (\emph{El baile}; \emph{La selva}.)
\end{quote}

\subsection*{What to know first}
\begin{itemize}
\raggedright
  \item el campo — the bigger version.
  \item la vaca — what stands in it.
\end{itemize}

\begin{sounds}
\begin{itemize}
\raggedright
  \item \emph{PRA-do}, two beats with the weight on the first. The \emph{pr} runs together, the \emph{r} is a single tap, and the \emph{d} between vowels is soft. $\to$ pronunciation reference
\end{itemize}
\end{sounds}

\begin{cousinweb}
\textbf{el prado} — \textquotedblleft{}meadow.\textquotedblright{}

Latin \textbf{pratum} was a meadow, and Spanish kept it almost unchanged.

English has the same word, gone a long way round. French built \textbf{prairie} on a Vulgar Latin \emph{prataria}, and English borrowed that — which is why the grasslands of North America are named in Latin by way of Paris.

The best use of this word is a place you may already know by name without knowing what it says. The \textbf{Prado} in Madrid is a museum, and the word on the front of it means \emph{meadow}. It stands on what were the \textbf{Prado de los Jerónimos}, the fields beside the monastery of Saint Jerome. People called the building \textquotedblleft{}the Prado\textquotedblright{} long before anybody made that official, and the Academy's dictionary still gives a second sense for \emph{prado}: a pleasant place in a town where people go to walk.

Where \emph{pratum} itself came from, nobody has established. The trail stops at Latin, which is a perfectly respectable place for a trail to stop.
\end{cousinweb}

\begin{grammarlens}[title={What you've built}]
You can now name a field of grass: \emph{el prado}.
\end{grammarlens}

\subsection*{Your turn}
Say these aloud:

\begin{itemize}
\raggedright
  \item \emph{el prado}
  \item \emph{en el prado}
  \item \emph{¿dónde está el prado?}
\end{itemize}

\begin{tcolorbox}[breakable,colback=teal!4,colframe=teal!35!black,title={Before you move on}]
Which English word for grassland comes from \emph{pratum}? (\emph{Prairie}, through French.) And what does the name of Madrid's Prado museum mean? (Meadow — it stands on one.)
\end{tcolorbox}
\section[el volcán]{el volcán — a mountain named after a god}

\label{lesson:ES-C378-volcan}



\begin{quote}
Say \emph{the jungle}, then \emph{the meadow}. (\emph{La selva}; \emph{El prado}.)
\end{quote}

\subsection*{What to know first}
\begin{itemize}
\raggedright
  \item la montaña — what it looks like from far off.
  \item el fuego — what comes out of it.
\end{itemize}

\begin{sounds}
\begin{itemize}
\raggedright
  \item \emph{vol-CÁN}, two beats with the weight on the second. The written accent is needed because a word ending in \emph{-n} would otherwise be stressed on the first beat. $\to$ pronunciation reference
\end{itemize}
\end{sounds}

\begin{cousinweb}
\textbf{el volcán} — \textquotedblleft{}volcano.\textquotedblright{}

Almost every word you have for a landscape describes what the thing is like. This one names a \textbf{person}.

\textbf{Vulcānus} was the Roman god of fire and the forge. The Academy's dictionary routes the word into Spanish through \textbf{Portuguese} \emph{volcão} — sailors' vocabulary, brought back — and from there to the god's name.

The Romans looked at the burning mountains of Sicily and read them as his workshop: the smoke was a chimney, the noise was hammering.

One thing to hold loosely. You will often read that the word comes specifically from the island of \textbf{Vulcano}, off Sicily, taken to be the chimney of the forge. It is a good story and it is not what the etymological dictionaries say — they go from the god to the word and leave the island out. Where a decorative detail turns up in travel writing but not in the dictionaries, keep it as decoration.
\end{cousinweb}

\begin{grammarlens}[title={What you've built}]
You can now name a mountain that burns: \emph{un volcán}.
\end{grammarlens}

\subsection*{Your turn}
Say these aloud:

\begin{itemize}
\raggedright
  \item \emph{el volcán}
  \item \emph{un volcán}
  \item \emph{¿dónde está el volcán?}
\end{itemize}

\begin{tcolorbox}[breakable,colback=teal!4,colframe=teal!35!black,title={Before you move on}]
Which god is \emph{volcán} named after? (Vulcan, the Roman god of fire and the forge.) And which language did Spanish take the word from? (Portuguese.)
\end{tcolorbox}
\section[la cumbre]{la cumbre — the ridge of a roof}

\label{lesson:ES-C378-cumbre}



\begin{quote}
Say \emph{the meadow}, then \emph{the volcano}. (\emph{El prado}; \emph{El volcán}.)
\end{quote}

\subsection*{What to know first}
\begin{itemize}
\raggedright
  \item la montaña — what it is the top of.
  \item el techo — what the word meant first.
\end{itemize}

\begin{sounds}
\begin{itemize}
\raggedright
  \item \emph{CUM-bre}, two beats with the weight on the first. The \emph{mb} runs together, and the \emph{r} is a single tap. $\to$ pronunciation reference
\end{itemize}
\end{sounds}

\begin{cousinweb}
\textbf{la cumbre} — \textquotedblleft{}summit.\textquotedblright{}

Latin \textbf{culmen} was the \textbf{ridge of a roof} — the beam along the very top, where the two slopes meet. The mountain came later.

Spanish still had the old sense well into the Middle Ages. Berceo, writing around 1236, uses \emph{cunbre} for a roofline, not a peak.

The ending is a pattern you already own without having been told its name. Latin nouns in \emph{-men} lose their middle vowel in Spanish and grow a \emph{b} to hold the wreckage together:

\begin{itemize}
\raggedright
  \item \emph{hominem} $\to$ \textbf{hombre};
  \item \emph{nominem} $\to$ \textbf{nombre};
  \item \emph{culminem} $\to$ \textbf{cumbre}.
\end{itemize}

Only the first of those three is real Latin. \emph{Nomen} and \emph{culmen} are neuter nouns, and a neuter accusative is the same as its nominative -- so Latin had \emph{nomen} and \emph{culmen}, full stop. \emph{Nominem} and \emph{culminem} are \textbf{reconstructed}: forms nobody ever wrote down, remade later by analogy with words like \emph{hominem}, which is real. Dictionaries print a star in front of a form like that to say exactly this.

The vowel, on the other hand, misbehaves. Latin \emph{culmen} had a \textbf{short u}, and a short u normally comes out as \emph{o} in Spanish — \emph{lupum} gives \emph{lobo}, \emph{buccam} gives \emph{boca}. By the ordinary route this word should be \emph{\textbf{combre}}. Corominas puts the blame on the \emph{l}, which pulled the vowel upward and then vanished into it, the same thing that happened in \emph{mucho} and \emph{escuchar}.

So do not use \emph{cumbre} as your model for what a short u does.

One last thing, because it is the opposite of what you would guess. The word closest to \emph{culmen} is not English \textbf{culminate}, which is a late borrowing taken straight off \emph{culmen} itself.

It is Latin \textbf{columen}. \emph{Culmen} is that word with its middle syllable squeezed out -- the two are one word in two states of compression.

Latin \textbf{columna}, the pillar, is built on the same root by a different route. That makes it a \textbf{cousin}, not a twin, and the difference is the whole point: a twin is one word in two shapes; a cousin is a separate word built from the same material. A roof ridge, a pillar and a mountaintop are all the thing that stands highest.
\end{cousinweb}

\begin{grammarlens}[title={What you've built}]
You can now name the top of a mountain: \emph{la cumbre}.
\end{grammarlens}

\subsection*{Your turn}
Say these aloud:

\begin{itemize}
\raggedright
  \item \emph{la cumbre}
  \item \emph{en la cumbre}
  \item \emph{¿dónde está la cumbre?}
\end{itemize}

\begin{tcolorbox}[breakable,colback=teal!4,colframe=teal!35!black,title={Before you move on}]
What did Latin \emph{culmen} mean first? (The ridge of a roof.) And which Latin word is \emph{culmen} the squeezed-down form of? (\emph{Columen} — while \emph{columna}, the pillar, is a cousin rather than a twin.)
\end{tcolorbox}
\section[el edificio]{el edificio — a house-making}

\label{lesson:ES-C378-edificio}



\begin{quote}
Say \emph{the volcano}, then \emph{the summit}. (\emph{El volcán}; \emph{La cumbre}.)
\end{quote}

\subsection*{What to know first}
\begin{itemize}
\raggedright
  \item la casa — the smaller version.
  \item la ciudad — what enough of them make.
\end{itemize}

\begin{sounds}
\begin{itemize}
\raggedright
  \item \emph{e-di-FI-cio}, four beats with the weight on the third. The \emph{c} before \emph{i} is soft, and the \emph{io} at the end is one glide. $\to$ pronunciation reference
\end{itemize}
\end{sounds}

\begin{cousinweb}
\textbf{el edificio} — \textquotedblleft{}building.\textquotedblright{}

\textbf{Edificio} is Latin \textbf{aedificium}, and it comes apart cleanly:

\begin{itemize}
\raggedright
  \item \textbf{aedes} — a temple in the singular, a house in the plural;
  \item \textbf{facere} — to make, which you have already met behind \emph{hacer}.
\end{itemize}

A building is a house-making. Spanish has kept both halves visible, which is unusual: most compounds this old have worn smooth.

English took the same Latin twice. An \textbf{edifice} is the thing built. To \textbf{edify} is the act — and English kept only the figurative sense of it, which is building a \emph{person} up rather than a wall. Spanish did the same: the Academy's dictionary gives \emph{edificar} a second meaning about instilling virtue.

One temptation to resist. \emph{Aedes} is sometimes glossed \textquotedblleft{}hearth,\textquotedblright{} and it is not a Latin meaning. It is a reconstructed prehistoric sense, inferred from a root meaning \textquotedblleft{}to burn\textquotedblright{} — a much older and much shakier claim than the ones above. In actual Latin, \emph{aedes} is a temple or a house, and that is what the word was built on.
\end{cousinweb}

\begin{grammarlens}[title={What you've built}]
You can now name what a street is made of: \emph{un edificio}.
\end{grammarlens}

\subsection*{Your turn}
Say these aloud:

\begin{itemize}
\raggedright
  \item \emph{el edificio}
  \item \emph{un edificio}
  \item \emph{¿dónde está el edificio?}
\end{itemize}

\begin{tcolorbox}[breakable,colback=teal!4,colframe=teal!35!black,title={Before you move on}]
What two Latin words make up \emph{aedificium}? (\emph{Aedes}, temple or house, and \emph{facere}, to make.) And what does English \emph{edify} build? (A person, not a wall.)
\end{tcolorbox}
